% =====================================================================
%  Lecture 16 -- Network Functions Virtualisation (NFV)    (Persian edition)
% =====================================================================
\documentclass[11pt,a4paper]{article}
\providecommand\ACNroot{../..}
\input{\ACNroot/templates/lecture-fa}

\ACNset{number}{16}
\ACNset{title}{مجازی‌سازی کارکردهای شبکه (\lr{NFV})}
\ACNset{shorttitle}{مجازی‌سازی کارکردهای شبکه}
\ACNset{titleother}{Network Functions Virtualisation (NFV)}
\ACNset{part}{بخش سوم -- معماری‌های نرم‌افزارمحور (درس‌های ۱۵ تا ۱۸)}
\ACNset{time}{۶ تا ۸ ساعت (۳ ساعت مطالعه، ۳ تا ۵ ساعت تمرین)}
\ACNset{level}{متوسط}
\ACNset{prereq}{درس‌های ۱۱ و ۱۵؛ سیستم‌عامل (فرایندها، حافظهٔ مجازی)؛ احتمال مقدماتی}
\ACNset{gen}{5}

\begin{document}
\ACNcover

% ---------------------------------------------------------------------
\section{اهداف یادگیری}
پس از مطالعهٔ این درس باید بتوانید:
\begin{objectives}
  \item \textbf{توضیح دهید} چرا اپراتورها کارکردهای شبکه را از دستگاه‌های اختصاصی به نرم‌افزار روی زیرساخت مشترک بردند.
  \item چارچوب معماری \lr{ETSI NFV} -- \acro{VNF}ها، \acro{NFVI} و کارکردهای \acro{MANO} با نقاط مرجعشان -- را \textbf{توصیف کنید}.
  \item ماشین‌های مجازی و کانتینرها را \textbf{مقایسه کنید} و کارکردهای شبکهٔ ابرزاد را \textbf{توضیح دهید}.
  \item چرخهٔ عمر \acro{VNF}، از جمله مقیاس‌دهی و ترمیم، را \textbf{دنبال کنید}.
  \item دسترس‌پذیری پشته‌های مجازی‌شده، نیاز مقیاس‌دهی و بودجهٔ پردازش هر بسته را \textbf{محاسبه کنید}.
  \item \acro{NFV} را به \acro{SDN} و به \acro{EPC} مجازی‌شده و هستهٔ ابرزاد نسل پنجم \textbf{مرتبط کنید}.
\end{objectives}

% ---------------------------------------------------------------------
\section{دانش پیش‌نیاز}
باید بدانید هستهٔ سیستم‌عامل، فرایند و حافظهٔ مجازی چیستند و دسترس‌پذیری سری و موازی چگونه ترکیب می‌شوند (درس ۱۱).
\acro{SDN} (درس ۱۵) همراه طبیعی \acro{NFV} است.

\paragraph{واژه‌های کلیدی} واژگان این درس همراه با معادل انگلیسی آن‌ها:
\begin{keyterms}
مجازی‌سازی کارکردهای شبکه & \lr{NFV} & اجرای کارکردهای شبکه به‌صورت نرم‌افزار روی کارسازهای همه‌منظوره. \\
کارکرد شبکهٔ مجازی / مؤلفه & \lr{VNF / VNFC} & یک کارکرد شبکهٔ مجازی‌شده / یکی از مؤلفه‌های آن. \\
زیرساخت \acro{NFV} & \lr{NFVI} & سخت‌افزار محاسبه، ذخیره‌سازی و شبکه و لایهٔ مجازی‌سازی. \\
مدیریت و هماهنگ‌سازی & \lr{MANO} & هماهنگ‌ساز \acro{NFV}، مدیر \acro{VNF} و مدیر زیرساخت مجازی‌شده. \\
توصیفگر \acro{VNF} / خدمت شبکه & \lr{VNFD / NSD} & الگوهایی که یک \acro{VNF} / یک خدمت شبکه را توصیف می‌کنند. \\
ابرناظر & \lr{Hypervisor} & نرم‌افزاری که ماشین‌های مجازی را روی یک میزبان اجرا می‌کند. \\
کانتینر & \lr{Container} & گروه فرایند ایزوله که هستهٔ میزبان را به اشتراک می‌گذارد. \\
کارکرد شبکهٔ ابرزاد & \lr{CNF} & کارکرد شبکه‌ای ساخته‌شده از کانتینرها/ریزخدمت‌ها. \\
فناوری‌های شتاب \lr{I/O} & \lr{SR-IOV / DPDK} & دسترسی مستقیم ماشین مجازی به کارت شبکه / پردازش بسته در فضای کاربر. \\
زنجیرهٔ کارکرد خدمت & \lr{Service function chain} & دنبالهٔ مرتب کارکردهایی که یک جریان از آن‌ها می‌گذرد. \\
\end{keyterms}

% ---------------------------------------------------------------------
\section{مقدمه و انگیزه}
تا دههٔ ۲۰۱۰ هر کارکرد شبکه -- \acro{MME}، \acro{HSS}، دروازه، دیوارهٔ آتش، متعادل‌کنندهٔ بار -- یک جعبهٔ اختصاصی بود:
سخت‌افزار ویژه، سیستم‌عامل اختصاصی و نرم‌افزار کارکرد که با هم فروخته می‌شدند. افزودن ظرفیت به معنای خرید و نصب جعبه‌های
بیشتر بود؛ معرفی یک خدمت به معنای یکپارچه‌سازی دستگاهی دیگر؛ و هر نوع، قطعات یدکی، ابزارها و چرخهٔ عمر خود را داشت. در
۲۰۱۲ گروهی از اپراتورها در یک «کاغذ سفید» که گروه مشخصات صنعتی \lr{ETSI NFV} را راه انداخت پیشنهاد کردند این کارکردها
به‌صورت نرم‌افزار روی کارسازهای استاندارد در مراکز داده و با هماهنگ‌سازی مشترک اجرا شوند \cite{etsinfv002}.

مزایای وعده‌داده‌شده همان مزایای رایانش ابری است: سخت‌افزار مشترک و هزینهٔ کمتر، ظرفیتی که تقاضا را دنبال می‌کند (مقیاس‌دهی
کشسان)، معرفی سریع‌تر خدمات و خودکارسازی. هزینه‌ها نیز واقعی‌اند: پردازش بسته روی پردازنده‌های همه‌منظوره دشوارتر از
تراشه‌های ویژه است و دسترس‌پذیری اکنون به پشته‌ای از لایه‌های سخت‌افزار، مجازی‌سازی و نرم‌افزار وابسته است که تیم‌های
مختلفی آن‌ها را اداره می‌کنند. \acro{EPC} مجازی‌شده نخستین کاربرد بزرگ موبایل بود؛ هستهٔ نسل پنجم از آغاز به‌صورت
نرم‌افزار ابرزاد طراحی شد (درس‌های ۱۸ و ۱۹).

% ---------------------------------------------------------------------
\section{مفاهیم بنیادی}

\subsection{چارچوب \lr{ETSI NFV}}
این چارچوب (\cref{fig:framework}) \cite{etsinfv002} سه حوزه دارد:
\begin{itemize}
  \item \textbf{\acro{VNF}ها}، هر یک با یک مدیر عنصر (\acro{EM}) برای پیکربندی سطح کاربرد.
  \item \textbf{\acro{NFVI}}: منابع فیزیکی محاسبه، ذخیره‌سازی و شبکه که یک لایهٔ مجازی‌سازی آن‌ها را به منابع مجازی انتزاع
        می‌کند.
  \item \textbf{\acro{MANO}} \cite{etsinfvman001}: \emph{هماهنگ‌ساز \acro{NFV}} (\acro{NFVO}) خدمات شبکه و منابع سراسری را
        مدیریت می‌کند؛ \emph{مدیر \acro{VNF}} (\acro{VNFM}) چرخهٔ عمر نمونه‌های \acro{VNF} را؛ و \emph{مدیر زیرساخت مجازی‌شده}
        (\acro{VIM}، مثلاً صفحهٔ کنترل \lr{OpenStack} یا \lr{Kubernetes}) منابع \acro{NFVI} را.
\end{itemize}
نقاط مرجع نام‌دار (\lr{Os-Ma, Or-Vnfm, Vi-Vnfm, Or-Vi, Ve-Vnfm, Nf-Vi, Vn-Nf, Vl-Ha}) تعیین می‌کنند چه کسی با چه کسی
گفت‌وگو می‌کند تا مؤلفه‌های فروشندگان مختلف ترکیب‌پذیر باشند.

\begin{acnfigure}
  \resizebox{\linewidth}{!}{\acnfig{lecture-16/nfv-framework}}
  \caption{چارچوب معماری \lr{ETSI NFV} با نقاط مرجع اصلی آن \cite{etsinfv002}.}
  \label{fig:framework}
\end{acnfigure}

\subsection{توصیفگرها و خدمات شبکه}
یک \acro{VNF} به‌صورت بسته‌ای شامل تصاویر نرم‌افزاری و یک \emph{توصیفگر \acro{VNF}} (\acro{VNFD}) تحویل می‌شود: مؤلفه‌های
آن (\acro{VNFC})، نیاز منابعشان، پیوندهای مجازی داخلی، نقاط اتصال، جنبه‌های مقیاس‌دهی و اسکریپت‌های چرخهٔ عمر. یک \emph{خدمت
شبکه} \acro{VNF}ها و پیوندهای مجازی میان آن‌ها را ترکیب می‌کند که با یک توصیفگر خدمت شبکه (\acro{NSD}) و، هرجا ترافیک باید
به ترتیب از کارکردها بگذرد، با یک گراف هدایت توصیف می‌شود. توصیفگرها استقرار را اعلانی می‌کنند: هماهنگ‌ساز یک الگو را به
نمونه‌های در حال اجرا تبدیل می‌کند.

\subsection{ماشین‌های مجازی و کانتینرها}
یک \emph{ماشین مجازی} یک سیستم‌عامل میهمان کامل را روی یک ابرناظر اجرا می‌کند؛ جداسازی قوی می‌دهد و از هر سیستم‌عاملی
پشتیبانی می‌کند، به بهای حافظه و زمان راه‌اندازی. یک \emph{کانتینر} گروهی از فرایندها را با فضاهای نام و گروه‌های کنترلی
هسته جدا می‌کند و هستهٔ میزبان را به اشتراک می‌گذارد؛ در چند ثانیه آغاز می‌شود و سربار اندکی دارد (\cref{fig:stacks}).
\emph{کارکردهای شبکهٔ ابرزاد} (\acro{CNF}) به‌صورت مجموعه‌هایی از ریزخدمت‌های کانتینری ساخته می‌شوند، حالت نشست را در
مخزن‌های دادهٔ جداگانه نگه می‌دارند تا نمونه‌ها دورریختنی باشند و با یک هماهنگ‌ساز کانتینر مانند \lr{Kubernetes} مستقر،
مقیاس‌دهی و ترمیم می‌شوند.

\begin{acnfigure}
  \resizebox{\linewidth}{!}{\acnfig{lecture-16/virt-stacks}}
  \caption{از دستگاه‌های اختصاصی تا ماشین‌های مجازی و کانتینرها.}
  \label{fig:stacks}
\end{acnfigure}

% ---------------------------------------------------------------------
\section{توضیحات فنی تفصیلی}

\subsection{مدیریت چرخهٔ عمر}
\acro{VNFM} یک \acro{VNF} را در عملیات چرخهٔ عمر مشخصات \acro{MANO} پیش می‌برد \cite{etsinfvman001}: \emph{ورود} (بسته
اعتبارسنجی و تصاویرش ذخیره می‌شود)، \emph{نمونه‌سازی}، \emph{مقیاس‌دهی} افقی (افزودن یا حذف نمونه‌های \acro{VNFC}) یا
عمودی (تغییر منابع آن‌ها)، \emph{ترمیم} (جایگزینی نمونه‌های خراب)، \emph{به‌روزرسانی/ارتقا} و \emph{خاتمه}. مقیاس‌دهی و ترمیم
را می‌تواند \acro{VNF} یا \acro{EM}، \acro{VNFM} بر اساس داده‌های کارایی و خطا، یا \acro{NFVO} برای کل یک خدمت شبکه آغاز
کند (\cref{fig:scale}).

\subsection{صفحه‌های دادهٔ پرکارایی}
دروازه‌ها و \acro{UPF}ها باید میلیون‌ها بسته در ثانیه را هدایت کنند. سوئیچینگ مجازی ساده از طریق هسته و ابرناظر پاسخ‌گو
نیست، بنابراین گره‌های \acro{NFVI} برای کارکردهای صفحهٔ کاربر از این‌ها استفاده می‌کنند:
\begin{itemize}
  \item \textbf{\acro{DPDK}}: راه‌انداز با حالت سرکشی و پردازش بسته در فضای کاربر، بدون وقفه و کپی‌های هسته؛
  \item \textbf{\lr{SR-IOV}}: کارت شبکه توابع مجازی‌ای عرضه می‌کند که مستقیماً به ماشین‌های مجازی یا کانتینرها داده می‌شوند؛
  \item \textbf{سنجاق‌کردن پردازنده و آگاهی از \acro{NUMA}}: هسته‌های اختصاصی روی سوکتی که به کارت شبکه وصل است و صفحه‌های
        بزرگ حافظه برای بافرهای بسته؛
  \item \textbf{تخلیه به سخت‌افزار}: کارت‌های شبکهٔ هوشمند یا سوئیچ‌های برنامه‌پذیر (درس ۱۵) برای کپسوله‌سازی، طبقه‌بندی یا
        رمزنگاری.
\end{itemize}
این روش‌ها بخشی از انعطاف \acro{NFV} (قابلیت حمل، تخصیص بیش از ظرفیت) را با کارایی معاوضه می‌کنند؛ بودجهٔ هر بسته در
\cref{sec:cycles} توضیح می‌دهد چرا لازم‌اند.

\subsection{تاب‌آوری در جهان مجازی}
در یک دستگاه اختصاصی، افزونگی درون جعبه ساخته می‌شد (بردهای افزونه، حالت همگام). در \acro{NFV} تاب‌آوری ویژگی طراحی
نرم‌افزار است: نمونه‌ها در میزبان‌ها و سایت‌ها پخش می‌شوند (قواعد ضدهم‌نشینی)، حالت تکثیر یا بیرونی می‌شود و نمونه‌های خراب
را هماهنگ‌ساز ترمیم می‌کند. خرابی‌های هم‌شکل در پشته بالا می‌روند: یک به‌روزرسانی معیوب ابرناظر یا اشکال هماهنگ‌ساز می‌تواند
همهٔ کارکردها را هم‌زمان متأثر کند.

\subsection{\lr{NFV} در شبکه‌های موبایل}
\acro{EPC} مجازی (\lr{vEPC}) \acro{MME}، \lr{S/P-GW}، \acro{HSS} و \acro{PCRF} را به‌صورت \acro{VNF} بسته‌بندی کرد، در آغاز
با «جابه‌جایی» نرم‌افزار دستگاه‌ها به ماشین‌های مجازی. هستهٔ نسل پنجم گام بعدی را برمی‌دارد: کارکردهای شبکه به‌صورت خدمت
تعریف می‌شوند (درس ۱۸)، به‌صورت \acro{CNF} پیاده و روی سکوهای مبتنی بر \lr{Kubernetes} مستقر می‌شوند و صفحهٔ کاربر در
سایت‌های لبه قرار می‌گیرد. در \acro{RAN}، \acro{RAN} مجازی و ابری باند پایه (\acro{CU} و \acro{DU}) را روی کارسازهایی با
سخت‌افزار شتاب‌دهنده اجرا می‌کنند.

% ---------------------------------------------------------------------
\section{مدل‌های ریاضی}

\subsection{دسترس‌پذیری یک پشتهٔ مجازی‌شده}\label{sec:avail}
\paragraph{مدل} یک نمونهٔ \acro{VNF} به سخت‌افزار میزبان، لایهٔ مجازی‌سازی و نرم‌افزار خود با دسترس‌پذیری‌های
$a_{\mathrm{hw}}$، $a_{\mathrm{v}}$ و $a_{\mathrm{sw}}$ (مستقل) نیاز دارد. به‌صورت سری، و با $k$ نمونهٔ افزونهٔ مستقل که یکی
کافی است،
\begin{equation}\label{eq:series}
  A_1 = a_{\mathrm{hw}}\,a_{\mathrm{v}}\,a_{\mathrm{sw}}, \qquad
  A_{k} = 1-(1-A_1)^{k}.
\end{equation}
\paragraph{مثال} $a_{\mathrm{hw}} = 0.999$، $a_{\mathrm{v}} = 0.9995$ و $a_{\mathrm{sw}} = 0.999$: $A_1 = 0.9975$، یعنی ۲۱٫۹ ساعت
قطعی در سال -- بدتر از هر مؤلفه. دو نمونه روی میزبان‌های مستقل: $A_2 = 1 - (0.0025)^2 = 0.999994$، حدود ۱۹۷ ثانیه در سال، به
شرط آنکه جایگزینی آنی باشد.
\paragraph{تفسیر} افزونگی در سطح نمونه بیش از جبران لایه‌های کم‌اطمینان‌تر را می‌کند، اما فقط اگر خرابی‌ها مستقل باشند
(میزبان‌ها، قفسه‌ها، سایت‌های متفاوت) و حالت از جایگزینی جان سالم به در برد.

\subsection{مقیاس‌دهی کشسان}\label{sec:scale}
\paragraph{مدل} هر نمونه $\mu$ درخواست در ثانیه را خدمت می‌دهد؛ اپراتور بهره‌وری را حداکثر $\rho_{\mathrm{t}}$ نگه می‌دارد. برای
بار ارائه‌شدهٔ $\lambda$ تعداد نمونه‌ها
\begin{equation}\label{eq:scale}
  n = \left\lceil \frac{\lambda}{\rho_{\mathrm{t}}\,\mu} \right\rceil
\end{equation}
است. اگر بار بتواند با نرخ $g$ (درخواست در ثانیه به ازای هر ثانیه) رشد کند و آماده‌شدن نمونهٔ جدید $T_\mathrm{i}$ طول بکشد،
ظرفیت آزاد باید رشد در طول نمونه‌سازی را پوشش دهد:
\begin{equation}\label{eq:headroom}
  n\mu - \lambda \ge g\,T_\mathrm{i}
  \quad\Longleftrightarrow\quad
  \rho \le 1 - \frac{g\,T_\mathrm{i}}{n\mu}.
\end{equation}
\paragraph{مثال} $\lambda = 40\,000$\,s$^{-1}$، $\mu = 5000$\,s$^{-1}$ و $\rho_{\mathrm{t}} = 0.7$: $n = \lceil 11.4\rceil = 12$. با
$g = 100$\,s$^{-2}$ و $T_\mathrm{i} = 60$ ثانیه، ظرفیت آزاد باید ۶۰۰۰\,s$^{-1}$ باشد، پس $\rho \le 1 - 6000/60\,000 = 0.9$: هدف ۷۰٪
امن است. کانتینرها ($T_\mathrm{i}$ چند ثانیه‌ای) بهره‌وری بسیار بیشتری از ماشین‌های مجازی (چند دقیقه‌ای) مجاز می‌کنند.

\subsection{بودجهٔ پردازش هر بسته}\label{sec:cycles}
\paragraph{مدل} با نرخ خط $R$ و بسته‌های $L$ بایتی (به‌علاوهٔ ۲۰ بایت پیش‌درآمد اترنت و فاصلهٔ میان قاب‌ها)، یک هسته با
بسامد ساعت $f$ این تعداد چرخه به ازای هر بسته دارد:
\begin{equation}\label{eq:cycles}
  c = \frac{f}{R/\bigl(8(L+20)\bigr)} .
\end{equation}
\paragraph{مثال} $f = 2.5$ گیگاهرتز، $R = 10$ گیگابیت بر ثانیه و $L = 64$: $14.88$ میلیون بسته در ثانیه و $c = 168$ چرخه --
در بدترین حالت کمتر از هزینهٔ یک خطای نهان‌گاه تا حافظهٔ اصلی. در ۱۰۰ گیگابیت بر ثانیه فقط ۱۷ چرخه باقی می‌ماند
(\cref{fig:cycles}).
\paragraph{تفسیر} هدایت بسته‌های کوچک با نرخ خط به هسته‌های موازی بسیار، دسته‌پردازی، ساختارهای دادهٔ سازگار با نهان‌گاه و
تخلیه نیاز دارد؛ ترافیک واقعی با بسته‌های میانگین بزرگ‌تر بودجه را به‌طور چشمگیری گشاده‌تر می‌کند.

\begin{acnfigure}
  \acnfig{lecture-16/cycles}
  \caption{چرخه‌های در دسترس به ازای هر بسته روی یک هستهٔ ۲٫۵ گیگاهرتزی با نرخ خط (\cref{eq:cycles}).}
  \label{fig:cycles}
\end{acnfigure}

\subsection{جایابی به‌صورت بسته‌بندی ظرف}\label{sec:pack}
\paragraph{مدل} \acro{VNFC}هایی با نیاز پردازشی $d_j$ باید روی کارسازهایی با ظرفیت $C$ قرار گیرند. تعداد کارسازها دست‌کم
\begin{equation}\label{eq:binpack}
  m \ge \left\lceil \frac{\sum_j d_j}{C} \right\rceil
\end{equation}
است و یافتن کمینه، مسئلهٔ \lr{NP}-سخت بسته‌بندی ظرف است؛ روش‌های ابتکاری مانند «اولین جای مناسب به ترتیب نزولی» به آن
نزدیک می‌شوند. ضدهم‌نشینی (نسخه‌ها روی کارسازهای متفاوت) و قیدهای تأخیر (صفحهٔ کاربر در لبه) قیدهایی می‌افزایند که
می‌توانند $m$ را افزایش دهند.

% ---------------------------------------------------------------------
\section{معماری شبکه و نمودارها}
یک استقرار \acro{NFV} اپراتوری معمولاً شامل چند مرکز دادهٔ مرکزی (کارکردهای صفحهٔ کنترل، پایگاه‌های دادهٔ مشترک،
\acro{OSS}) و سایت‌های لبهٔ بسیار (کارکردهای صفحهٔ کاربر، باند پایهٔ \acro{RAN}) است که همه با یک لایهٔ هماهنگ‌سازی مشترک
مدیریت می‌شوند. کنترل‌گرهای \acro{SDN} (درس ۱۵) تارهای مراکز داده و انتقال میان سایت‌ها را پیکربندی می‌کنند تا پیوندهای مجازی
خدمات شبکه از ابتدا تا انتها برقرار شوند. \Cref{fig:framework} نمای کارکردی را نشان می‌دهد؛ از نظر فیزیکی، خود مؤلفه‌های
\acro{MANO} به‌صورت کاربردهای مجازی‌شده روی همان زیرساخت اجرا می‌شوند.

% ---------------------------------------------------------------------
\section{رویه‌های گام‌به‌گام}

\begin{procedurebox}{نمونه‌سازی یک خدمت شبکه}\label{proc:inst}
\begin{steps}
  \item اپراتور (\acro{OSS}) نمونه‌سازی یک خدمت شبکه را با ارجاع به یک \acro{NSD} واردشده از \acro{NFVO} درخواست می‌کند.
  \item \acro{NFVO} منابع را بررسی، \acro{VIM}ها/سایت‌ها را انتخاب و از \acro{VNFM} هر \acro{VNF} نمونه‌سازی آن را درخواست می‌کند.
  \item \acro{VNFM} از \acro{NFVO} مجوز منابع می‌گیرد، سپس از \acro{VIM} می‌خواهد ماشین‌های مجازی یا کانتینرها، شبکه‌ها و
        ذخیره‌سازی را مطابق \acro{VNFD} بسازد.
  \item \acro{VIM} منابع را روی \acro{NFVI} تخصیص می‌دهد و تصاویر را راه‌اندازی می‌کند؛ \acro{VNFM} نمونه‌ها را پیکربندی و به
        \acro{EM} اطلاع می‌دهد.
  \item \acro{NFVO} پیوندهای مجازی میان \acro{VNF}ها را (احتمالاً از طریق یک کنترل‌گر \acro{SDN}) می‌سازد و خدمت را نمونه‌سازی‌شده
        گزارش می‌دهد.
\end{steps}
\end{procedurebox}

\begin{procedurebox}{مقیاس‌دهی افقی خودکار}\label{proc:scale}
\begin{steps}
  \item \acro{VNF} یا \acro{EM} باری بالاتر از آستانهٔ تعریف‌شده در \acro{VNFD} گزارش می‌کند (یا \acro{VNFM} آن را از سنجه‌ها
        تشخیص می‌دهد).
  \item \acro{VNFM} گام مقیاس‌دهی را تعیین و از \acro{NFVO} مجوز می‌گیرد.
  \item \acro{VIM} منابع را تخصیص و یک نمونهٔ جدید \acro{VNFC} را آغاز می‌کند.
  \item \acro{VNFM} آن را پیکربندی و به توزیع بار داخلی \acro{VNF} می‌افزاید؛ ترافیک به سوی آن جابه‌جا می‌شود.
\end{steps}
\end{procedurebox}

\begin{acnfigure}
  \resizebox{\linewidth}{!}{\acnfig{lecture-16/scale-out}}
  \caption{مقیاس‌دهی افقی سادهٔ \acro{VNF} (\cref{proc:scale}).}
  \label{fig:scale}
\end{acnfigure}

% ---------------------------------------------------------------------
\section{مثال‌های حل‌شده}

\begin{examplebox}{طراحی افزونگی}\label{ex:avail}
\textbf{مسئله.} با دسترس‌پذیری مؤلفه‌های \cref{sec:avail}، یک نمونه، دو نمونه روی یک میزبان و دو نمونه روی میزبان‌های متفاوت
را مقایسه کنید.

\textbf{حل.} یک نمونه: $A_1 = 0.9975$ (۲۱٫۹ ساعت در سال). یک میزبان: سخت‌افزار و ابرناظر مشترک‌اند، پس
$A = a_{\mathrm{hw}}a_{\mathrm{v}}\,[1-(1-a_{\mathrm{sw}})^2] = 0.9985$ (۱۳٫۱ ساعت در سال): فقط خرابی‌های نرم‌افزاری پوشش داده
می‌شوند. میزبان‌های متفاوت: $A_2 = 0.999994$ (۱۹۷ ثانیه در سال). ضدهم‌نشینی ضروری است.
\end{examplebox}

\begin{examplebox}{مقیاس‌دهی یک کارکرد سیگنالینگ}\label{ex:scale}
\textbf{مسئله.} یک کارکرد شبیه \acro{AMF} باید ساعت اوج $40\,000$ رویه در ثانیه را پردازش کند؛ هر نمونه ۵۰۰۰ در ثانیه را. با
بهره‌وری هدف ۷۰٪ چند نمونه لازم است و پس از خرابی یک نمونه بهره‌وری چقدر است؟

\textbf{حل.} $n = 12$ (\cref{eq:scale})؛ بهره‌وری $40\,000/60\,000 = 67\,\%$. پس از یک خرابی: $40\,000/55\,000 = 73\,\%$ -- هنوز زیر
اشباع، و هماهنگ‌ساز در این فاصله نمونه را ترمیم می‌کند.
\end{examplebox}

\begin{examplebox}{بسته‌بندی \lr{VNFC}ها}\label{ex:pack}
\textbf{مسئله.} \acro{VNFC}ها به ۶، ۵، ۵، ۴، ۴، ۳، ۳ و ۲ هسته نیاز دارند؛ کارسازها ۱۶ هسته دارند. کران پایین و یک جایابی «اولین
جای مناسب به ترتیب نزولی» بیابید.

\textbf{حل.} $\sum d_j = 32$، پس $m \ge 2$. اولین جای مناسب نزولی: کارساز ۱ مقادیر ۶، ۵، ۵ (۱۶ هسته)؛ کارساز ۲ مقادیر ۴، ۴، ۳، ۳،
۲ (۱۶ هسته) را می‌گیرد. دو کارساز، که بهینه است -- اما بدون ظرفیت آزاد برای ترمیم.
\end{examplebox}

% ---------------------------------------------------------------------
\section{مقایسهٔ فناوری‌ها}
\begin{table}[htbp]
\centering\small
\caption{دستگاه‌های فیزیکی، \acro{VNF}های مبتنی بر ماشین مجازی و \acro{CNF}های ابرزاد.}\label{tab:cmp}
\rowcolors{2}{acnmist}{white}
\begin{tabularx}{\linewidth}{@{}l Y Y Y@{}}
\toprule
\rowcolor{white} & دستگاه اختصاصی & \acro{VNF} (ماشین مجازی) & \acro{CNF} (کانتینر) \\
\midrule
سخت‌افزار        & اختصاصی، ویژهٔ فروشنده   & کارسازهای تجاری مشترک & کارسازهای تجاری مشترک \\
جداسازی          & فیزیکی                    & ابرناظر               & فضاهای نام هسته، \lr{cgroups} \\
زمان راه‌اندازی   & ساعت‌ها تا روزها (نصب)   & چند دقیقه             & چند ثانیه \\
مقیاس‌دهی        & خرید و نصب جعبه           & افزودن ماشین مجازی (افقی/عمودی) & افزودن خودکار پاد \\
رسیدگی به حالت   & درون جعبه                 & اغلب درون ماشین مجازی  & مخزن‌های دادهٔ بیرونی \\
چرخهٔ عمر        & ابزارهای ویژهٔ فروشنده    & \acro{MANO} (\lr{NFVO, VNFM, VIM}) & \lr{Kubernetes} + اپراتورها، \lr{CI/CD} \\
سرعت صفحهٔ داده  & نرخ خط تراشهٔ ویژه        & نیاز به \acro{DPDK}، \lr{SR-IOV} & نیاز به \acro{DPDK}، \lr{SR-IOV}، تخلیه \\
\bottomrule
\end{tabularx}
\end{table}

\paragraph{\lr{NFV} و \lr{SDN}} این دو مکمل و مستقل‌اند: \acro{NFV} تغییر می‌دهد کارکردها \emph{کجا} اجرا شوند (نرم‌افزار روی
کارسازهای مشترک)، \acro{SDN} تغییر می‌دهد هدایت \emph{چگونه} کنترل شود (دستگاه‌های برنامه‌پذیر، کنترل مرکزی). با هم اجازه
می‌دهند یک خدمت شبکه -- کارکردها و پیوندهای میان آن‌ها -- با نرم‌افزار ایجاد شود.

% ---------------------------------------------------------------------
\section{نکته‌های کلیدی}
\begin{takeawaybox}
\begin{itemize}
  \item \acro{NFV} کارکردهای شبکه را به‌صورت نرم‌افزار روی زیرساخت مشترک اجرا می‌کند که \acro{MANO} (\lr{NFVO, VNFM, VIM}) آن
        را مدیریت می‌کند.
  \item توصیفگرها (\lr{VNFD, NSD}) استقرار را اعلانی می‌کنند؛ عملیات چرخهٔ عمر شامل مقیاس‌دهی و ترمیم است.
  \item کانتینرها و طراحی ابرزاد (نمونه‌های بی‌حالت، حالت بیرونی) پایهٔ هستهٔ نسل پنجم‌اند.
  \item دسترس‌پذیری سری با هر لایه کاهش می‌یابد؛ نمونه‌های افزونهٔ مستقل آن را بازمی‌گردانند: $A_k = 1-(1-A_1)^k$.
  \item مقیاس‌دهی: $n = \lceil\lambda/(\rho_{\mathrm{t}}\mu)\rceil$ با ظرفیت آزاد $gT_\mathrm{i}$؛ بودجهٔ بسته $c = 8f(L+20)/R$ چرخه.
\end{itemize}
\end{takeawaybox}

% ---------------------------------------------------------------------
\section{تمرین‌ها}

\begin{exercisebox}{1}{چه کسی چه می‌کند؟}\label{exr:mano}
کدام مؤلفهٔ \acro{MANO} (\acro{NFVO}، \acro{VNFM}، \acro{VIM}) مسئول این‌هاست: (الف) ساختن یک ماشین مجازی؛ (ب) تصمیم به افزودن یک
نمونهٔ \acro{VNFC}؛ (ج) اتصال چند \acro{VNF} در یک خدمت شبکه؛ (د) نگه‌داشتن سیاههٔ منابع \acro{NFVI}؟
\end{exercisebox}

\begin{exercisebox}{1}{ماشین مجازی یا کانتینر؟}\label{exr:vmct}
یک استدلال برای اجرای \acro{HSS} مجازی‌شده در ماشین‌های مجازی و یکی برای اجرای آن به‌صورت کانتینر بیاورید.
\end{exercisebox}

\begin{exercisebox}{2}{افزونگی با دامنهٔ مشترک}\label{exr:host}
با \cref{eq:series} زمان قطعی سالانهٔ سه نمونه روی سه میزبان متفاوت را محاسبه کنید که یک دامنهٔ برق و شبکهٔ مشترک با
دسترس‌پذیری ۰٫۹۹۹۹ دارند. کدام جمله غالب است؟
\end{exercisebox}

\begin{exercisebox}{2}{بودجهٔ بسته}\label{exr:cycles}
یک هستهٔ \acro{UPF} با ۳ گیگاهرتز کار می‌کند و به ۴۰۰ چرخه برای هر بسته نیاز دارد. یک هسته با بسته‌های ۶۴ و ۸۰۰ بایتی چه
گذردهی‌ای را تحمل می‌کند؟ برای ۱۰۰ گیگابیت بر ثانیه با ۸۰۰ بایت چند هسته لازم است؟
\end{exercisebox}

\begin{exercisebox}{3}{تأخیر مقیاس‌دهی}\label{exr:lag}
بار با $g = 200$\,s$^{-2}$ افزایش می‌یابد. نمونه‌ها $\mu = 4000$\,s$^{-1}$ را خدمت می‌دهند. بیشینهٔ بهره‌وری امن را برای $n = 10$
نمونه وقتی آماده‌شدن نمونهٔ جدید ۱۲۰ ثانیه (ماشین مجازی) یا ۱۰ ثانیه (کانتینر) طول می‌کشد مقایسه کنید. چه چیز دیگری سرعت
مقیاس‌دهی را محدود می‌کند؟
\end{exercisebox}

\begin{exercisebox}{3}{جایابی با ضدهم‌نشینی}\label{exr:pack}
\acro{VNFC}های \cref{ex:pack} شامل دو نسخه از یک مؤلفهٔ ۶ هسته‌ای (\lr{A})، دو نسخه از یک مؤلفهٔ ۵ هسته‌ای (\lr{B}) و بقیه
(۴، ۴، ۳، ۳، ۲) بدون قیدند؛ نسخه‌های یک مؤلفه باید روی کارسازهای متفاوت باشند و هر کارساز باید ۲ هسته برای ترمیم آزاد نگه
دارد. جایابی‌ای با کمترین تعداد کارساز ۱۶ هسته‌ای بیابید.
\end{exercisebox}

\subsection*{پاسخ‌ها}
\addcontentsline{toc}{subsection}{پاسخ‌ها}

\begin{solutionbox}{\ref{exr:mano}}
(الف) \acro{VIM}؛ (ب) \acro{VNFM} (احتمالاً به درخواست \lr{EM/VNF} یا \acro{NFVO})؛ (ج) \acro{NFVO}؛ (د) \acro{VIM} (با دید
سراسری‌ای که \acro{NFVO} نگه می‌دارد).
\end{solutionbox}

\begin{solutionbox}{\ref{exr:vmct}}
ماشین مجازی: جداسازی قوی برای یک پایگاه دادهٔ حساس امنیتی که کلیدهای مشترکین را نگه می‌دارد و پشتیبانی از تصویر موجود
فروشنده. کانتینر: مقیاس‌دهی و ترمیم سریع، ردپای کوچک‌تر و یکپارچگی با ابزارهای ابرزاد، اگر نرم‌افزار \acro{HSS} با حالت
بیرونی و تکثیرشده طراحی شده باشد.
\end{solutionbox}

\begin{solutionbox}{\ref{exr:host}}
سه نمونهٔ مستقل: $1-(0.0025)^3 \approx 1-1.6\times10^{-8}$. دامنهٔ مشترک به‌صورت سری: $A = 0.9999\times(1-1.6\times10^{-8})
\approx 0.99990$، یعنی حدود ۵۳ دقیقه در سال که تقریباً همه ناشی از دامنهٔ مشترک است. استقلال دامنه‌های خرابی از تعداد نمونه‌ها
مهم‌تر است.
\end{solutionbox}

\begin{solutionbox}{\ref{exr:cycles}}
یک هسته: $3\times10^9/400 = 7.5$ میلیون بسته در ثانیه. بسته‌های ۶۴ بایتی: $7.5\times10^6\times84\times8 = 5.0$ گیگابیت بر
ثانیه؛ بسته‌های ۸۰۰ بایتی: $7.5\times10^6\times820\times8 = 49$ گیگابیت بر ثانیه (روی سیم). برای ۱۰۰ گیگابیت بر ثانیه با ۸۰۰
بایت: $100/49 = 2.03$، یعنی ۳ هسته (با ظرفیت آزاد بیشتر).
\end{solutionbox}

\begin{solutionbox}{\ref{exr:lag}}
ظرفیت $n\mu = 40\,000$\,s$^{-1}$. ماشین مجازی: ظرفیت آزاد $200\times120 = 24\,000$، پس $\rho \le 0.4$؛ کانتینر:
$200\times10 = 2000$، پس $\rho \le 0.95$. محدودیت‌های دیگر: تأخیر تشخیص حلقهٔ پایش، بارگیری تصویر، گرم‌شدن (نهان‌گاه‌ها،
همگام‌سازی حالت)، پیکربندی مجدد متعادل‌کنندهٔ بار و در دسترس بودن سخت‌افزار آزاد در سایت.
\end{solutionbox}

\begin{solutionbox}{\ref{exr:pack}}
ظرفیت قابل استفاده ۱۴ هسته برای هر کارساز است؛ نیاز کل ۳۲، پس دست‌کم ۳ کارساز. جایابی: کارساز ۱: \lr{A} (۶) + \lr{B} (۵) +
۳ = ۱۴؛ کارساز ۲: \lr{A} (۶) + \lr{B} (۵) + ۳ = ۱۴؛ کارساز ۳: ۴ + ۴ + ۲ = ۱۰. ضدهم‌نشینی برقرار است (هر کارساز حداکثر یک نسخه از
\lr{A} و \lr{B} دارد) و روی هر کدام ۲ هسته آزاد می‌ماند: سه کارساز.
\end{solutionbox}

% ---------------------------------------------------------------------
\section{مطالعهٔ بیشتر}
چارچوب معماری \acro{NFV} در \lr{ETSI GS NFV 002} \cite{etsinfv002} و مدیریت و هماهنگ‌سازی در \lr{ETSI GS NFV-MAN 001}
\cite{etsinfvman001} آمده است. کروتز و همکاران \cite{kreutz2015sdn} رابطهٔ \acro{SDN} و \acro{NFV} را بررسی می‌کنند؛ رومر و
همکاران \cite{rommer2020core} هستهٔ ابرزاد نسل پنجم را توصیف می‌کنند؛ کلاینراک \cite{kleinrock1975queueing} مدل‌های صف به‌کاررفته
در اندازه‌گذاری را پوشش می‌دهد.

\begin{latin}
\printbibliography[heading=subbibliography,title={References}]
\end{latin}
\end{document}
